\documentclass[11pt]{article}
\usepackage[T1]{fontenc}
\usepackage[utf8]{inputenc}
\usepackage{amsmath,amssymb}
\usepackage{booktabs}
\usepackage[margin=1in]{geometry}

\title{Evaluation Metrics}
\date{}

\begin{document}

\maketitle

To evaluate the performance of the trained models, we use several quantitative and qualitative metrics. These metrics measure overall predictive performance, class-specific behavior, and training dynamics.

\section{Confusion Matrix}

The primary evaluation tool is the confusion matrix.

For a classification problem with $N$ classes, the confusion matrix is an $N \times N$ table where:

\begin{itemize}
	\item rows represent the true class labels,
	\item columns represent the predicted class labels.
\end{itemize}

Each entry $CM_{i,j}$ indicates the number of samples whose true label is class $i$ but were predicted as class $j$. Correct predictions lie on the main diagonal of the matrix, while off-diagonal entries correspond to misclassifications.

From the confusion matrix, we can define several quantities for a specific class $i$:

\begin{itemize}
	\item \textbf{True Positive ($\mathrm{TP}_i$):} samples correctly predicted as class $i$.
	\item \textbf{False Positive ($\mathrm{FP}_i$):} samples predicted as class $i$ but belonging to another class.
	\item \textbf{False Negative ($\mathrm{FN}_i$):} samples belonging to class $i$ but predicted as another class.
	\item \textbf{True Negative ($\mathrm{TN}_i$):} samples that are neither predicted nor labeled as class $i$.
\end{itemize}

These values allow us to compute several evaluation metrics.

\section{Overall Accuracy}

The overall accuracy measures the proportion of correctly classified samples:

$$
\mathrm{Accuracy} = \frac{\text{number of correct predictions}}{\text{total number of samples}}
$$

Using the confusion matrix $CM$:

$$
\mathrm{Accuracy} = \frac{\mathrm{tr}(CM)}{\sum CM}
$$

where $\mathrm{tr}(CM)$ is the sum of the diagonal entries. Accuracy provides a high-level measure of model performance.

\section{Class-Specific Reliability}

Because per-class accuracy is mathematically inflated by high True Negative rates in multi-class problems, we rely on precision, recall, and the F1-score to rigorously evaluate the model's strength on individual classes.

\begin{itemize}
	\item \textbf{Precision ($\mathrm{Pre}_i$):} Measures prediction reliability. Of all samples predicted as class $i$, how many were actually class $i$?
\end{itemize}

$$
\mathrm{Pre}_i = \frac{\mathrm{TP}_i}{\mathrm{TP}_i + \mathrm{FP}_i}
$$

\begin{itemize}
	\item \textbf{Recall ($\mathrm{Rec}_i$):} Measures detection capability. Of all actual samples of class $i$, how many did the model correctly identify?
\end{itemize}

$$
\mathrm{Rec}_i = \frac{\mathrm{TP}_i}{\mathrm{TP}_i + \mathrm{FN}_i}
$$

\begin{itemize}
	\item \textbf{F1-Score ($\mathrm{F1}_i$):} The harmonic mean of precision and recall, providing a single balanced metric for class $i$ performance.
\end{itemize}

$$
\mathrm{F1}_i = 2 \frac{\mathrm{Pre}_i \cdot \mathrm{Rec}_i}{\mathrm{Pre}_i + \mathrm{Rec}_i}
$$

We synthesize these into a comprehensive metrics table, calculating the macro-average, the unweighted mean of the per-class metrics, to ensure minority classes are weighted equally in the final evaluation.

\section{Macro-Averaged Metrics}

To summarize performance across all classes, we compute macro averages, which treat each class equally:

$$
\mathrm{Metric}_{\mathrm{macro}} = \frac{1}{N} \sum_{i=1}^{N} \mathrm{Metric}_i
$$

Macro-averaged precision, recall, and F1 provide a single metric describing performance across the entire dataset.

\section{Metrics Summary Table}

All class-level metrics are summarized in the following format:

\begin{center}
\begin{tabular}{lccc}
\toprule
Class & Precision & Recall & F1 Score \\
\midrule
C1 &  &  &  \\
C2 &  &  &  \\
$\cdots$ &  &  &  \\
CN &  &  &  \\
\textbf{Macro Avg} &  &  &  \\
\bottomrule
\end{tabular}
\end{center}

This table allows us to identify which classes the model handles well and which classes are more difficult.

\section{Learning Curves and Generalization}

To diagnose underfitting or overfitting during training, we plot learning curves mapping both accuracy and loss against epochs for both the training and validation sets.

\begin{itemize}
	\item \textbf{Healthy Learning:} Both training and test metrics improve synchronously.
	\item \textbf{Overfitting:} Training loss continues to decrease while test loss stagnates or increases.
	\item \textbf{Underfitting:} Both training and test metrics stall at suboptimal levels.
\end{itemize}

We quantify the model's ability to extrapolate to unseen data using the generalization gap. While often calculated as $(\text{Training Accuracy} - \text{Test Accuracy})$, the gap between training and test loss serves as a more sensitive early-warning indicator of overfitting.

\section{Qualitative Error Analysis}

Statistical metrics identify where the model fails; qualitative analysis determines why. By isolating and reviewing misclassified samples, false positives and false negatives, we categorize errors into systematic vulnerabilities:

\begin{itemize}
	\item \textbf{Ambiguous Digits:} Cases where the input is objectively difficult to discern, for example a poorly drawn ``4'' resembling a ``9''.
	\item \textbf{Feature Mismatch:} Consistent confusion between structurally similar classes.
\end{itemize}

\section{Writer Bias}

Finally, we measure writer bias to evaluate whether the model performs unevenly across different handwriting styles. Let $A_i$ denote the accuracy on samples written by writer $i$. The bias is defined as:

$$
\mathrm{WriterBias} = \max(A_i) - \min(A_i)
$$

A large value indicates that the model performs significantly better on some writers than others.

\end{document}
